
% =====================================================================
\section{The setting: a bounded resolvable space}
\label{sec:setting}
% =====================================================================

We work throughout in a metric--measure space. The objects of the theory are
regions (subsets), their separators (boundaries), and partitions; the single
hypothesis constrains the ambient space to be bounded and finitely resolvable.

\subsection{Ambient space}

\begin{definition}[Resolvable space]
\label{def:space}
A \emph{resolvable space} is a triple $(\Space,d,\mu)$ where $(\Space,d)$ is a
metric space and $\mu$ is a Borel measure on $\Space$ that is positive on
non-empty open sets and finite on bounded sets. We write $\Leb$ for Lebesgue
measure when $\Space\subseteq\R^{m}$, and $\Haus^{s}$ for $s$-dimensional
Hausdorff measure~\citep{federer1969,mattila1995}.
\end{definition}

\begin{definition}[Region, complement, separator]
\label{def:region}
A \emph{region} is a Borel set $A\subseteq\Space$ with
$\mu(A)>0$ and $\mu(\Space\setminus A)>0$ (so both $A$ and its complement are
non-negligible). Its \emph{topological boundary} is
$\Bdy A=\cl(A)\cap\cl(\Space\setminus A)$. A \emph{separator} of $A$ is any
Borel set $\Sigma$ with $\Space\setminus\Sigma=U_A\sqcup U_{A^c}$ a disjoint
union of two non-empty relatively open sets with $A\setminus\Sigma\subseteq
U_A$ and $(\Space\setminus A)\setminus\Sigma\subseteq U_{A^c}$. A separator is
\emph{sharp} if $\mu(\Sigma)=0$.
\end{definition}

\begin{remark}
$\Bdy A$ is the smallest closed separator in the topological sense; a general
separator may be thicker. The question the paper turns on is whether a
\emph{sharp} separator (one of zero measure) can exist for the regions that an
embedded agent actually uses. The answer, under \Cref{ax:brs}, is no.
\end{remark}

\subsection{Partitions and resolution}

\begin{definition}[Finite measurable partition]
\label{def:partition}
A \emph{finite measurable partition} of $\Space$ is a finite family
$\Part=\{\cell_1,\dots,\cell_N\}$ of pairwise-disjoint Borel sets with
$\bigcup_i\cell_i=\Space$ and $\mu(\cell_i)>0$ for each $i$. A partition
$\Part'$ \emph{refines} $\Part$, written $\Part'\preceq\Part$, if every cell of
$\Part'$ is contained in a cell of $\Part$.
\end{definition}

\begin{definition}[Resolution]
\label{def:resolution}
A resolvable space \emph{carries resolution $\reso>0$} for a partition $\Part$
if every cell $\cell\in\Part$ satisfies $\diam(\cell)\ge\reso$; equivalently, no
two points closer than $\reso$ are forced into distinct cells. The resolution
is the finest scale at which the partition distinguishes points.
\end{definition}

\begin{remark}[Resolution is an agent-relative datum]
\label{rem:reso-relative}
Resolution is not a property of $\Space$ alone but of the partition an agent can
realise. Two agents resolving the same space may carry different $\reso$. The
floor we derive will inherit this relativity: its magnitude depends on $\reso$,
but its strict positivity will not.
\end{remark}

\subsection{The hypothesis}

\begin{axiom}[Bounded resolvable space]
\label{ax:brs}
The space $(\Space,d,\mu)$ satisfies:
\begin{enumerate}[label=\textup{(BRS\arabic*)},leftmargin=3.2em]
\item\label{brs:bdd} \textbf{Boundedness.} $\Space$ is bounded and
$0<\mu(\Space)<\infty$; moreover $\cl(\Space)$ is compact.
\item\label{brs:res} \textbf{Finite resolution.} Every partition realisable by
an embedded agent carries some resolution $\reso>0$; there is no partition into
cells of arbitrarily small diameter at fixed, finite descriptive cost.
\item\label{brs:hier} \textbf{Hierarchical refinability.} There is a sequence of
finite measurable partitions $\Part_0\succeq\Part_1\succeq\Part_2\succeq\cdots$,
each refining the last, generating the Borel $\sigma$-algebra modulo $\mu$-null
sets; each $\Part_n$ carries a resolution $\reso_n>0$.
\end{enumerate}
\end{axiom}

\begin{remark}[Status of the hypothesis]
\label{rem:brs-status}
\ref{brs:bdd} is the substantive geometric content: the accessible space is
finite in extent and measure, and its closure is compact (so the
extreme-value and finite-subcover arguments below apply). \ref{brs:res} is the
finitude of any embedded agent: an agent with finite descriptive resources
cannot specify cells of vanishing diameter---this is what forbids the limiting
zero-thickness cut. \ref{brs:hier} states that finer partitions exist and
exhaust the measurable structure, but each \emph{individual} partition is still
finite-resolution; refinement is unbounded in the limit but bounded at every
finite stage. The three clauses are exactly what is needed for the
Boundary--Thickness Theorem and nothing more.
\end{remark}

\begin{lemma}[Finiteness of cell count]
\label{lem:finite-cells}
Under \ref{brs:bdd} and \ref{brs:res}, every realisable partition $\Part$ has
finitely many cells, and the count is bounded by
$N(\Part)\le \mu(\Space)/\mu_{\min}(\Part)$, where
$\mu_{\min}(\Part)=\min_i\mu(\cell_i)>0$.
\end{lemma}

\begin{proof}
By \ref{brs:res} each cell has $\diam(\cell_i)\ge\reso$, and by
\ref{brs:bdd} $\mu(\cell_i)>0$ with the cells disjoint and of finite total
measure $\mu(\Space)<\infty$. If the count were infinite, then
$\sum_i\mu(\cell_i)=\mu(\Space)<\infty$ would force
$\inf_i\mu(\cell_i)=0$, contradicting that finitely-described cells carry a
uniform positive lower measure bound under finite resolution. Hence the count
is finite, and $N(\Part)\,\mu_{\min}(\Part)\le\sum_i\mu(\cell_i)=\mu(\Space)$
gives the bound.
\end{proof}
